% Project AllMight, Seven-Book edition, Book 01, Chapter 01 project unit.
% Intended as a chapter module; includes neither a preamble nor documentclass.
% C++23 companion: tools/labs/ines_header_cpp23_lab.cpp.
\section{A Cartridge Is a Contract, Not Just a Bag of Bytes}
\label{sec:allmight-nes01}

Imagine opening a game file. A naive program sees a sequence of bytes, but the
emulator must decide which bytes describe the file itself and which bytes belong
to a program or to graphics. The first sixteen bytes of a conventional iNES
1.0 file are a \emph{header}: an agreement about how to interpret the rest of
the file. Our first NES milestone is deliberately smaller than an emulator.
We will build a reader that proves the layout before any CPU tries to execute
an instruction. This is the foundation of trustworthy systems software.

\subsection{A bit does not explain itself}

One byte is eight binary digits. Under the ordinary unsigned positional rule,
its value is $\sum_{i=0}^{7} b_i2^i$, with each $b_i\in\{0,1\}$. This
explains why a byte can represent an integer from $0$ to $255$, but it does
not explain whether the byte is a count, a set of flags, a fragment of an
address, or text. A format supplies that missing meaning. To test bit $k$,
we form the mask $2^k$ and compute $x\mathbin{\&}2^k$; the result is nonzero
exactly when bit $k$ is set. To extract the high nibble, compute $x\gg4$.
These operations do not depend on whether a human writes the value in binary
or hexadecimal notation.

\subsection{Derive the cartridge layout}

An iNES 1.0 header begins with the four bytes
\texttt{4E 45 53 1A} (\texttt{NES} followed by \texttt{0x1A}).
Header byte 4 counts program ROM banks, each exactly $16\,384$ bytes;
byte 5 counts character ROM banks, each exactly $8\,192$ bytes.
Header byte 6 contains mirroring, battery, trainer, and four-screen flags
in its lower four bits, and the low four mapper-ID bits in its upper four.
Header byte 7 supplies the remaining high mapper-ID bits. Thus, for
header bytes $h_i$,
\[
\mathrm{mapper}=(h_6\gg4)\;\mathbin{|}\;(h_7\mathbin{\&}\mathrm{0xF0}).
\]
If byte 6 has its trainer flag set, a 512-byte trainer follows the 16-byte
header. The payload offsets are therefore
\[
\begin{aligned}
O_{\mathrm{PRG}} &= 16 + (\mathrm{trainer}\ ?\ 512:0),\\
S_{\mathrm{PRG}} &= 16384h_4,\\
O_{\mathrm{CHR}} &= O_{\mathrm{PRG}} + S_{\mathrm{PRG}},\\
S_{\mathrm{CHR}} &= 8192h_5.
\end{aligned}
\]
The conditional expression adds $512$ if the trainer flag is set, otherwise
zero. A zero CHR-ROM bank count usually indicates CHR RAM allocated by the
emulator; it does \emph{not} ask us to read 8 KiB
of nonexistent CHR-ROM bytes from the file. These rules describe legacy
iNES 1.0; NES 2.0 has different size semantics and must be parsed separately.

\subsection{Validate before constructing runtime state}

It would be tempting to declare a C++ \texttt{struct Header} and read file
bytes directly into it. That conflates a specified on-disk schema with a
compiler's object-layout rules and risks using partially read data. Instead,
our C++23 function receives a \texttt{std::span<const std::uint8\_t>},
checks that the header exists, checks the signature and version, calculates
and bounds-checks the payload regions, and returns
\texttt{std::expected<Layout, ParseError>}.
The success type names all validated offsets and flags; the error type
identifies the failure. Checking remaining length by subtraction after
checking the starting offset avoids the usual unsafe pattern of comparing
against a possibly overflowing end offset. No dynamic mapper object is
created while the data is still untrusted.

The unit is implemented in \texttt{tools/labs/ines\_header\_cpp23\_lab.cpp}.
It uses eleven synthetic cases that require no commercial game files. They
cover ordinary PRG/CHR banks, trainer displacement, mapper nibble extraction,
CHR RAM, invalid signatures, short headers, truncated trainer/PRG/CHR data,
an explicitly rejected NES 2.0 header, zero PRG banks, and four-screen flags.
Run it with a standards-mode C++23 compiler and warnings enabled. A passing
test establishes these declared cases only: it does not establish complete
iNES compatibility or a functioning NES emulator.

\subsection{Connect this milestone to the actual emulator}

The source reference is the author's
\texttt{NESEmulator/Cartridge.cpp} in the \texttt{NESEmulator} repository.
It already shows the intended cartridge-to-mapper architecture, but its
constructor reads a header before checking the read status, assumes payload
reads succeed, and can mark a cartridge valid without constructing a
supported mapper. Study that implementation as a \emph{diagnostic case},
not as an unquestionable answer. Later units will integrate validated
cartridge storage with a mapper interface, CPU bus, 6502 instruction tracing,
PPU registers, APU timing, and deterministic frame tests.

\subsection{Check your understanding}

Without running code, take a header with one PRG bank, two CHR banks, and a
trainer. Calculate the start and length of both payloads, then explain what
happens if the final CHR byte is missing. Modify the synthetic test fixture
to verify your prediction. Next, explain why the CPU address \texttt{0x1801}
refers to the same physical RAM byte as \texttt{0x0001} when the 2-KiB RAM is
mirrored across the low 8-KiB CPU address window. The solution uses a mask
\texttt{address \& 0x07FF}; it is a hardware mapping rule, not a general
property of all addresses.

\paragraph{Completion gate.} A learner can explain the header byte-by-byte,
derive the offsets, reject malformed images without undefined behavior,
and reproduce all eleven test cases. Neither a completed emulator nor a
validated production cartridge loader is claimed at this milestone.

\paragraph{Reference specifications.} Consult the NESdev iNES file-format
specification (www.nesdev.org/wiki/INES), the NES 2.0 format description
(www.nesdev.org/wiki/NES\_2.0), and NESdev's mirroring reference
(www.nesdev.org/wiki/Mirroring). The present lab deliberately rejects a
header that advertises NES 2.0, whose size and mapper fields differ. Old
legacy dumps can contain misleading reserved bytes; a production loader
requires an explicit compatibility policy rather than guessing.
